\section*{Chapitre \ref{ch:dv:stochastic-volatility} --- Volatilité stochastique}

\begin{solution}{exo:dv:stochastic-volatility:1}
Base~: $2\times1.5\times0.04-0.6^2=-0.24$~; calibré~: $2\times2.05\times0.060-0.66^2=-0.19$. Les deux échouent~: la
variance peut atteindre zéro, ce que gère la fonction caractéristique, mais une simulation naïve doit tronquer.
\end{solution}

\begin{solution}{exo:dv:stochastic-volatility:2}
$v(t)=\bar v+(v_0-\bar v)e^{-\kappa t}$~: 0.0589 à six mois, 0.0489 à un an. La variance totale à un an est
$0.04+0.04(1-e^{-1.5})/1.5=0.0607$~: une volatilité de 24.6~\%.
\end{solution}

\begin{solution}{exo:dv:stochastic-volatility:3}
$\varphi(-i)=\E[e^{x_T}]=\E[S_T]/F=1$ car le forward est l'espérance de $S_T$ sous la mesure de valorisation. Une
valeur de 1.01 signifie que le forward du modèle s'écarte de 1~\%~: chaque call est mal évalué d'environ 1~\% du forward, ce qui constitue un
arbitrage contre le marché des contrats à terme, signe le plus souvent d'une mauvaise branche ou d'une erreur de codage.
\end{solution}

\begin{solution}{exo:dv:stochastic-volatility:4}
Sous la mesure de valorisation, la dérive de la variance contient déjà la prime de risque~: la calibration sur options estime
$\kappa\bar v$ et $\kappa$ tels qu'ils sont valorisés, et non tels qu'ils se comportent. La volatilité historique
estime le niveau réel, typiquement inférieur à celui qui est valorisé (la prime de risque de variance du chapitre~14)~; fixer $\bar v$ à partir de
l'historique sous-évalue les options.
\end{solution}

\begin{solution}{exo:dv:stochastic-volatility:5}
Avec $\rho=0$, le prix est une moyenne de prix de Black--Scholes à forward fixé~; en \omterm{def:dv:implied-volatility-and-its-surface:surface}{log-moneyness}, la
volatilité implicite de Black--Scholes d'une telle moyenne est symétrique autour de $k=0$ (put et call en $\pm k$ se correspondent par
symétrie put--call). Au voisinage de la monnaie, le prix de Black--Scholes est presque linéaire en volatilité, si bien que le prix moyen
correspond à la volatilité moyenne $\E[\bar\sigma]$, qui est inférieure à $\sqrt{\E[\bar\sigma^2]}$, selon
l'inégalité de Jensen.
\end{solution}

\begin{solution}{exo:dv:stochastic-volatility:6}
Le skew court est piloté par $\rho$ et $\eta$, le skew long aussi, et $\kappa$ contrôle la vitesse à laquelle le skew décroît avec
l'échéance~: n'accentuer que l'extrémité courte exige à la fois un $\kappa$ plus élevé et un $\eta$ plus grand, ce qui modifie aussi la
courbure et la structure par terme. Celle du skew de Heston a une forme fixe, et un événement local dans une partie de
celle-ci ne peut être ajustée en déplaçant un seul paramètre.
\end{solution}

\begin{solution}{exo:dv:stochastic-volatility:7}
$\eta=0.77$~; le skew à un an 90--110 monte à 6.45 points, puisque, avec une corrélation négative, la \omterm{def:dv:stochastic-volatility:sv}{volatilité de
la volatilité} accroît aussi la pente du skew.
\end{solution}

\begin{solution}{exo:dv:stochastic-volatility:8}
Le modèle a été calibré là où il s'ajuste~; à une semaine, son skew vaut la moitié de celui du marché, si bien qu'il sous-évalue les
puts à court terme hors de la monnaie et donne de mauvaises couvertures dans cette zone. Valoriser les options hebdomadaires à partir du smile à court terme du marché, ou
avec un modèle qui produit des skews courts pentus (sauts, volatilité rugueuse).
\end{solution}

\begin{solution}{pb:dv:stochastic-volatility:1}
\textbf{1.} 23.4~\%, 20.1~\%, 16.9~\%, 14.4~\%, 13.3~\%.
\textbf{2.} Skew de 5.71 points~; courbure de 0.30 point.
\textbf{3.} Non~: $2\kappa\bar v-\eta^2=-0.24$.
\textbf{4.} Moyenner sur une variance aléatoire abaisse la volatilité moyenne (Jensen), et la corrélation négative
abaisse encore la volatilité à la monnaie~: 16.9~\% à un an.
\textbf{5.} 17.4~\%, encore inférieur à 20~\% pour la même raison.
\textbf{6.} 3.94 points.
\textbf{7.} $-1.77$ points.
\textbf{8.} Peu~: la corrélation déplace la probabilité entre les queues, le forward étant fixé, si bien que le
prix à la monnaie ne varie que d'un montant du second ordre.
\textbf{9.} Les risk reversals, et tout produit dont la valeur dépend du skew~: barrières, digitales, autocallables.
\textbf{10.} Cette volatilité tend à monter quand l'indice baisse. Une baisse de 1~\% est un choc
$dW^1=-0.01/\sqrt{v}=-0.05$~; comme $d\sqrt v\approx\tfrac12\eta\,dW^2$ et $\E[dW^2\mid dW^1]=\rho\,dW^1$, la
hausse espérée de $\sqrt v$ vaut $\tfrac12\times0.6\times0.7\times0.05=0.0105$~: environ un point de volatilité.
\textbf{11.} $\eta=0.77$.
\textbf{12.} Le skew monte jusqu'à 6.45 points~: avec une corrélation négative, la \omterm{def:dv:stochastic-volatility:sv}{volatilité de la volatilité} ajoute
du skew en plus de la courbure.
\textbf{13.} Rendre la corrélation moins négative jusqu'à ce que le skew retrouve sa cible.
\textbf{14.} $2\times1.5\times0.04-0.77^2=-0.47$.
\textbf{15.} Les options sur les ailes, la convexité des swaps de variance (chapitre~14), les cliquets et les produits sur la \omterm{def:dv:stochastic-volatility:sv}{volatilité de la volatilité}.
\textbf{16.} Les paramètres interagissent (ici $\eta$ déplace aussi le skew)~; les déplacer un à un laisse les autres
faux.
\textbf{17.} $\kappa$ et $\eta$ sont presque redondants pour une surface allant jusqu'à quelques années~; fixer $\kappa$ supprime une
direction plate de l'ajustement et rend les autres paramètres stables d'un jour à l'autre (chapitre~24).
\textbf{18.} Son skew plafonne là où celui du marché continue de s'accentuer~: la moitié de celui du marché à une semaine.
\textbf{19.} $-1.77$ points de volatilité pour $+0.2$ en $\rho$~; $\eta=0.77$ pour une courbure de 0.5 point.
\textbf{20.} L'incertitude sur la variance future, qui fait monter le prix des options dont la valeur est convexe en
volatilité~: les ailes.
\end{solution}

\begin{solution}{iq:dv:stochastic-volatility:1}
$dS=\sqrt vS\,dW^1$, $dv=\kappa(\bar v-v)dt+\eta\sqrt v\,dW^2$, $d\langle W^1,W^2\rangle=\rho\,dt$. $v_0$~: le niveau
à court terme~; $\bar v$~: le niveau à long terme~; $\kappa$~: la vitesse à laquelle l'un devient l'autre, et celle à laquelle le skew s'estompe~; $\eta$~: la
courbure du smile (ailes)~; $\rho$~: son skew.
\iqlookfor{une phrase par paramètre, sur la surface.}
\end{solution}

\begin{solution}{iq:dv:stochastic-volatility:2}
Deux mouvements browniens, un seul actif risqué négocié~: le risque de variance ne peut pas être couvert avec l'action. Couvrir le delta avec
l'action et le \omterm{def:dv:greeks-and-the-hedging-pnl:greeks}{vega} (exposition à la variance) par une autre option, typiquement de même échéance~; la seconde option rend
le \omterm{def:dv:no-arbitrage-and-the-fundamental-theorems:complete}{marché complet} dans le modèle.
\iqlookfor{le décompte des risques face aux instruments.}
\end{solution}

\begin{solution}{iq:dv:stochastic-volatility:3}
Conditionner sur la trajectoire entière de la variance~: à corrélation nulle, $\ln S_T$ est normal de variance $\int v\,dt$, si bien que le
prix conditionnel est celui de Black--Scholes à cette variance~; on moyenne sur les trajectoires de variance (la \omterm{prop:dv:stochastic-volatility:mixing}{formule de mélange}).
\iqlookfor{conditionnement et indépendance.}
\end{solution}

\begin{solution}{iq:dv:stochastic-volatility:4}
Le logarithme complexe de la fonction caractéristique saute de branche (le «~petit piège de Heston~»)~: utiliser la forme avec
$g=(\xi-d)/(\xi+d)$ et $e^{-dT}$, qui reste sur la branche principale~; vérifier aussi le domaine d'intégration et
l'amortissement de l'intégrale de Fourier.
\iqlookfor{coupures de branche, puis numérique.}
\end{solution}

\begin{solution}{iq:dv:stochastic-volatility:5}
Vérifier d'abord les entrées (une mauvaise cotation, une échéance périmée)~; vérifier ensuite si l'ajustement a une direction plate ($\rho$ et $\eta$
qui se compensent)~; stabiliser la calibration (fixer $\kappa$, pénaliser les variations de paramètres, chapitre~24) et mesurer l'effet sur le P\&L
du saut sur les produits exotiques avant de l'accepter~; si le marché a vraiment bougé, les réserves sur les produits exotiques doivent le refléter.
\iqlookfor{données, identifiabilité, stabilité, impact sur le P\&L.}
\end{solution}

\begin{solution}{iq:dv:stochastic-volatility:6}
Son smile à court terme ne provient que de la diffusion de la variance sur un court intervalle, si bien que le skew à la monnaie tend vers une
constante proportionnelle à $\rho\eta/\sqrt{v_0}$ lorsque l'échéance diminue, alors que celui du marché croît. Remèdes~: ajouter des sauts
(chapitre~13), qui produisent du skew aux courtes échéances, ou rendre la volatilité rugueuse (chapitre~12), ce qui produit un skew
qui explose comme une puissance de l'échéance.
\iqlookfor{la limite aux échéances courtes, et les deux familles de remèdes.}
\end{solution}
